% U3 WS3 -- Type A/B: solutions, molarity, particulate views. Block 6.
\documentclass{apchem-worksheet}

\apcunit{3}
\apcunittitle{Properties of Substances and Mixtures}
\apcdoctitle{Worksheet 3 \textbullet\ Solutions and Concentration}
\apcsubtitle{CED 3.7, 3.8 \textbullet\ ionic solutes are drawn as separated ions, always}

\begin{document}
\wsheader[When a problem involves an ionic solute, state how many moles of
each ion form before computing concentrations.]

\problem[3.7]{Molarity calculations:}
\begin{parts}
  \item \SI{0.400}{\mole} in \SI{500.}{\milli\liter}:
        \answerblank[2.8cm]{\SI{0.800}{\Molar}}
  \item \SI{25.0}{\gram} \ce{NaOH} in \SI{2.00}{\liter}:
        \answerblank[2.8cm]{\SI{0.313}{\Molar}}
  \item Moles in \SI{150.}{\milli\liter} of \SI{0.250}{\Molar}:
        \answerblank[2.8cm]{\SI{0.0375}{\mole}}
  \item Mass of \ce{KNO3} in \SI{250.}{\milli\liter} of
        \SI{0.100}{\Molar}: \answerblank[2.8cm]{\SI{2.53}{\gram}}
\end{parts}

\problem[3.7]{Dilutions:}
\begin{parts}
  \item \SI{25.0}{\milli\liter} of \SI{12.0}{\Molar} diluted to
        \SI{500.}{\milli\liter}: \answerblank[2.8cm]{\SI{0.600}{\Molar}}
  \item Volume of \SI{6.00}{\Molar} needed to make
        \SI{250.}{\milli\liter} of \SI{1.50}{\Molar}:
        \answerblank[2.8cm]{\SI{62.5}{\milli\liter}}
  \item Explain why $M_1V_1 = M_2V_2$ works, in one sentence.
        \answerlines[2]{Adding solvent changes the volume but not the number
        of moles of solute, so the product $MV$ (which equals moles) is
        conserved.}
\end{parts}

\problem[3.7]{Ion concentrations:}
\begin{parts}
  \item $[\ce{Na+}]$ in \SI{0.20}{\Molar} \ce{Na3PO4}:
        \answerblank[2.8cm]{\SI{0.60}{\Molar}}
  \item $[\ce{SO4^2-}]$ in \SI{0.15}{\Molar} \ce{Al2(SO4)3}:
        \answerblank[2.8cm]{\SI{0.45}{\Molar}}
  \item Total ion concentration in \SI{0.10}{\Molar} \ce{CaCl2}:
        \answerblank[2.8cm]{\SI{0.30}{\Molar}}
\end{parts}

\problem[3.8]{In the box, draw a particulate representation of
\SI{0.10}{\Molar} \ce{MgBr2}. Use filled circles for \ce{Mg^2+} and open
circles for \ce{Br-}, show at least 3 magnesium ions, and omit water
molecules (note that you have done so). \workspace[4.4cm]
\keyonly{\begin{apcanswer}At least 3 filled and 6 open circles ($1{:}2$
ratio), all separate and randomly distributed --- never drawn as bonded
\ce{MgBr2} units. Caption: ``water omitted for clarity.''\end{apcanswer}}}

\problem[3.8]{A student draws \SI{0.1}{\Molar} \ce{KCl} as eight intact
``\ce{KCl}'' pairs floating in water. Identify the error and its
consequence. \answerlines[2]{Ionic compounds dissociate completely in
water; the drawing must show 8 separate \ce{K+} and 8 separate \ce{Cl-}.
The error predicts wrong conductivity and wrong particle counts.}}

\problem[3.7]{Describe how to prepare \SI{500.}{\milli\liter} of
\SI{0.200}{\Molar} \ce{CuSO4} from the solid, including the glassware.
\answerlines[3]{$n = 0.100$ mol $\times \SI{159.62}{\gram\per\mole} =
\SI{15.96}{\gram}$. Weigh it, dissolve in roughly 300 mL of water, transfer
quantitatively to a \SI{500}{\milli\liter} volumetric flask, then add water
to the calibration mark and invert to mix.}}

\begin{spiralstrip}{Unit 3 \textbullet\ blocks 1--5}
  \item Partial pressure of \ce{N2} if $\chi = 0.80$, $P_{\text{tot}} =
        \SI{2.5}{\atm}$: \answerblank[2.6cm]{\SI{2.0}{\atm}}
  \item Volume of \SI{0.50}{\mole} at STP: \answerblank[2.6cm]{\SI{11.2}{\liter}}
  \item Faster at equal $T$: \ce{O2} or \ce{H2}?
        \answerblank[2cm]{\ce{H2}}
  \item Strongest IMF in \ce{HF}: \answerblank[3.5cm]{hydrogen bonding}
  \item Real gases deviate most at: \answerblank[4cm]{low $T$, high $P$}
\end{spiralstrip}

\end{document}
